\section{Long-Term Vision}

\migaki{} should not become another agent framework. Its long-term ambition is to establish execution optimization as an independent systems layer.

Every successful computing platform eventually separated describing computation from executing computation. Programming languages did this with compilers. SQL did this with query planners. Distributed systems did this with schedulers. Cloud infrastructure did this with runtimes, orchestration, and containers. AI systems are approaching the same transition: applications and agent frameworks should describe intent and strategy, while an independent execution layer determines how that work is realized under explicit constraints.

In mature AI systems:

\begin{center}
\fbox{\begin{minipage}{0.86\linewidth}
Applications describe intent.\\
Agent frameworks determine strategy.\\
Execution optimizers determine realization.\\
Workflow engines preserve durable state.\\
Gateways provide provider connectivity.\\
Observability systems record behavior.\\
Evaluation systems measure outcomes.\\
Models perform computation.
\end{minipage}}
\end{center}

\migaki{} focuses on one responsibility:

\begin{quote}
Given a logical AI workflow or observed agent trajectory and explicit constraints, produce the best measurable execution plan or reuse decision and evidence for why it is acceptable.
\end{quote}

The long-term system should close the loop:

\begin{center}
Observed trajectories $\rightarrow$ reusable evidence graphs $\rightarrow$ optimized \mir{} plans $\rightarrow$ capability-aware execution $\rightarrow$ new evidence.
\end{center}

This loop is what turns observability into optimization. Traces alone explain the past. Plans alone describe intent. \migaki{} should connect the two so repeated agent work becomes progressively more reusable, more auditable, and more efficient without hiding the assumptions that made reuse possible.
